% ch-a5-01.tex —— 第 5 篇（Ben-Zvi–Chen–Helm–Nadler）篇首 + §1 + §2 前半
% 底本 PDF 434–445 页 = 印刷页 421–432
\chapter{凝聚与可构造局部 Langlands 对应之间}
\ysrc{421}
\begin{center}
{\small\itshape David Ben-Zvi \quad Harrison Chen \quad David Helm \quad David Nadler}
\end{center}
\medskip
\noindent\textbf{摘要.} 局部 Langlands 对应的各种精细形式，力求把局部域上约化群的表示与 Langlands 参数叠上的层联系起来。但究竟是哪一类层？Kazhdan--Lusztig 理论风格的猜想（源于 Vogan 与 Soergel）用可构造层描述一个群及其纯内型在固定中心特征下的表示。几何 Langlands 风格的猜想（源于 Fargues、Zhu 与 Hellmann）则用凝聚层描述与同晶（isocrystal）相关联的一大族群在中心特征变动下的表示。后一类猜想还发生在更大的参数空间上，其中允许 Frobenius（或复共轭）带幂幺部分。

本文提出一个在这两种设定之间插值的普遍机制。该机制源自循环同调理论，并借助导出代数几何中的圆作用来解读。我们把这一视角应用于阿基米德与非阿基米德局部域两类局部 Langlands 猜想的范畴形式。在非阿基米德情形，我们描述圆作用如何联系仿射 Hecke 代数以及 $\mathrm{GL}_n$ 全部光滑表示的凝聚实现与可构造实现，并就一般情形提出一个联系两种设定的机制。在阿基米德情形，我们说明如何用圆作用，从一个凝聚形式——Fargues 猜想的实类比，即扭量线上的驯顺分歧几何 Langlands 猜想（本文综述的对象）——导出可构造局部 Langlands 对应（Adams--Barbasch--Vogan 与 Soergel 所给出的形式）。
\footnote{2020 年数学主题分类：主分类 20G25。}
\footnote{受 NSF 个人基金 DMS-2001398、DMS-2302356（DBZ）以及 DMS-2101466、DMS-2401178（DN）资助。\textcopyright 2025 American Mathematical Society。}

\medskip
\noindent\textbf{目录.}
\begin{list}{}{\leftmargin=2.2em\labelwidth=2em\labelsep=0.4em\itemindent=-2.2em\itemsep=0.1em}
\item 1.\quad 概述：Langlands 参数上的层
\item 2.\quad 圈空间上的等变层
\item 3.\quad 非阿基米德局部 Langlands 与圆作用
\item 4.\quad 从阿基米德局部 Langlands 到扭量几何 Langlands
\item 附录 A.\quad 基础
\item 参考文献
\end{list}

\ysrc{422}
\section{概述：Langlands 参数上的层}

表示论中谱分解的基本主题，是要用对偶对象 $\widehat G$ 来描述群 $G$ 的表示；$\widehat G$ 参数化（适当的）不可约表示。调和分析要求把 $G$ 的大表示（例如 $G$-空间上的函数）谱式地描述为 $\widehat G$ 上的一族向量空间（重数空间）。我们还可能希望用不可约表示来描述不可分解模或标准模、解决扩张问题；或者更一般地，用 $\widehat G$ 的几何来描述表示的族。

这一主题是几何表示论中很大一部分内容的底层思想：人们力求把表示描述为参数空间上的层，并利用层的几何来解决表示论问题。起这种作用的层主要有两类：可构造层（取广义，含反常层与 $D$-模）与凝聚层（非正式地含拟凝聚层与 ind-凝聚层）。

注意：本文把注意力完全限制在特征零的表示论，并在导出（$\infty$-范畴）而非阿贝尔的框架下工作；因此所有范畴都是域 $k\supset\Q$ 上 $k$-线性的 dg 范畴（或 $k$-线性稳定 $\infty$-范畴）。

\subsection{表示论中的可构造层}

让我们回顾可构造层起作用的两个经典例子——Kazhdan--Lusztig 理论与 Springer 理论；它们分别是阿基米德域与非阿基米德域上局部 Langlands 对应的原型。在这些情形以及几何表示论中其他令人感兴趣的情形里，我们所处理的是只有有限多个轨道的群作用的等变性，因此可以等价地使用等变导出范畴的 Betti 版本或 de Rham 版本，也就是使用可构造层或 $D$-模。

Kazhdan--Lusztig 理论把复 Lie 代数 $\mathfrak g$ 的（无穷小特征平凡的）最高权模范畴 $\mathcal O_0$ 等同于旗簇 $\mathcal B=G/B$ 上 $N$-等变反常层的范畴。这一描述在范畴 $\mathcal O_0$ 的不可约对象与旗簇上 $N$-轨道的集合之间建立双射（后者单连通，因而只带唯一的局部系统）。更重要的是，它使人能够用旗簇的等变几何——由其层范畴所刻画——来描述不可分解模的合成列（这正是 Kazhdan--Lusztig 猜想的主题）。

在 Springer 理论中，我们考虑 Chevalley 群 $G(\F_q)$ 的复表示，特别是那些含有 $B(\F_q)$-不动向量的表示，即幂幺主序列表示。Springer 理论把这些表示实现于幂零锥 $\mathcal N\subset\mathfrak g$ 上 $G$-等变反常层范畴的内部。具体如下：幂幺主序列表示由 Morita 理论等同于有限 Hecke 代数
\[
\mathcal H^f = \C[B(\F_q)\backslash G(\F_q)/B(\F_q)] = \End_{G(\F_q)}(\C[G(\F_q)/B(\F_q)])
\]
的模；而 $\mathcal H^f$（在取定 $q$ 的一个平方根之后）可以等同于 $G$ 的 Weyl 群 $W$ 的群代数（不过这一等同对我们的故事而言在某种意义上只是附带的）。\ysrc{423}设 $\mu:\widetilde{\mathcal N}=T^*\mathcal B\to\mathcal N$ 记 Springer 解消，并考虑 Springer 层
\[
\mathbf S=\mu_*\C_{\widetilde{\mathcal N}/G}[\dim(G/B)]\in\operatorname{Perv}(\mathcal N/G).
\]
有限 Hecke 代数与 Springer 层的自同态之间有同构
\[
\mathcal H^f\simeq\End_{\mathcal N/G}(\mathbf S).
\]
Springer 层是投射的（事实上整个范畴 $\operatorname{Perv}(\mathcal N/G)$ 半单），于是我们得到 Hecke 模到反常层的一个满嵌入
\[
\{\text{$G(\F_q)$ 的幂幺主序列表示}\}\simeq\mathcal H^f\text{-mod}\xrightarrow{\ \sim\ }\langle\mathbf S\rangle\subset\operatorname{Perv}(\mathcal N/G),
\]
即嵌入为由 Springer 层生成的子范畴。于是，不可约幂幺主序列表示由幂零轨道连同其上的某些等变局部系统来分类；等价地说，由某个幂零元中心化子的分支群的表示来分类。（缺失的那些局部系统由 Lusztig 的广义 Springer 对应来补足。）

\subsection{局部 Langlands 纲领中的可构造层}

\ysrc{424}Vogan 表述下的局部 Langlands 对应 [Vo93] 提出了局部域上约化群表示的一个类似分类。固定局部域 $K$ 与特征零的域 $k$。为简单起见，再固定 $K$ 上连通分裂约化群 $G$，其 Langlands 对偶群为 $\check G$——$k$ 上的代数群。局部 Langlands 对应给出一个映射
\[
\{\text{$G(K)$ 的不可约光滑表示}\}\longrightarrow\mathcal L_{\check G}(K):=\{\mathcal L_K\longrightarrow\check G\ \text{连续}\}/\check G
\]
它把表示映为 Langlands 参数，即 $K$ 的 Weil（--Deligne）群到 $G$ 中的［连续］表示\tn{底本 423 页此处作「into $G$」，按上下文应为「into $\check G$」：Langlands 参数取在对偶群 $\check G$ 中。}。该映射的纤维，即 \emph{L-包}，预期由相应 Langlands 参数中心化子的分支群的某些不可约表示来参数化——也就是由某些不可约等变局部系统参数化。为看到那些缺失的局部系统，人们把单个群 $G(K)$ 换成它的纯内型的集合（当以层论方式思考 $G(K)$ 的表示时，这些纯内型自然出现；例如见 [Ber16] 与下文第 4.1 节）。于是我们得到一个猜想性的双射：Langlands 参数轨道上的等变局部系统，与 $G(K)$ 的纯内型的某一族不可约表示之间。这导向 Vogan 的猜想 [Vo93]：在 Grothendieck 群的水平上，用适当 Langlands 参数空间上的可构造层，完整描述 $G(K)$ 及其纯内型在固定中心特征下的表示论——其精神承自 Kazhdan--Lusztig 理论与 Springer 理论。相应的 Langlands 参数空间带有固定的无穷小特征（如 [Vo93] 所述）；特别地，这意味着在非阿基米德情形固定 $\check G$ 中的一个半单共轭类（某个半单 Frobenius 元的类），在阿基米德情形固定 $\check{\mathfrak g}$ 中的一个半单共轭类（来自 Weil 群作用的导数）。

在阿基米德情形，这一图景由 Adams--Barbasch--Vogan [ABV92] 建立（参看 Mason--Brown 的讲义 [MB25]）。此外 Soergel [So01] 猜想这些结果有一个范畴化加强，断言：实约化群的纯内型的 Harish-Chandra 模（固定无穷小特征）的范畴，与 ABV 空间上 Langlands 参数的可构造层的等变导出范畴 Koszul 对偶。Soergel 猜想在 [So01] 中对复数群与 $\mathrm{SL}_2(\R)$ 得到证明，对拟分裂块则在 [BV21] 中证明。

在非阿基米德情形，可构造层起中心作用的对应是 Kazhdan--Lusztig 理论 [KL87] 与 Ginzburg [CG97] 所建立的仿射类比，它用幂幺 Langlands 参数空间上的（某些）等变可构造层来刻画仿射 Hecke 代数（等价地，$G(K)$ 的非分歧主序列表示）的表示论。Lusztig [Lu95b]（承 [Lu88,Lu95a]）把这一图景补全为「仿射广义 Springer 对应」，在（$G$ 为伴随、单、非分歧时）$G(K)$ 的纯内型的全部幂幺表示与幂幺 Langlands 参数上的全部等变可构造层之间建立局部 Langlands 对应。这项工作近来由 Solleveld 推广：[Sol18,Sol19] 先推广到连通 $G$ 的全部幂幺表示，随后 [Sol22a,Sol22b] 证明了 Vogan 的 $p$-进 Kazhdan--Lusztig 猜想，覆盖了广泛得多的一类表示。

所有这些进展的主要输入，是把与表示的各块（block）相伴的 Hecke 代数等同于适当反常层的 Ext-代数。再结合 [Ri13,RiRu16,RiRu21,Sol22b] 中的形式性定理，就得到表示的范畴与 Langlands 参数上可构造层的范畴之间的导出等价。这一可构造图景的一个缺点是：它似乎不足以描述连续参数的变化——事实上，Vogan 猜想中出现的 Langlands 参数空间，作为无穷小特征的函数是间断变化的。

\subsection{局部 Langlands 纲领中的凝聚层}

凝聚层之所以起作用，其根源在于：交换代数 $Z$ 的全部模可以等同于 $\Spec(Z)$ 上的拟凝聚层。若 $Z$ 来自结合代数 $A$ 的中心，则得到从 $A$-模到 $\Spec(Z)$ 上拟凝聚层的局部化函子。更抽象地说，范畴 $\mathcal C$ 层化于其 Bernstein 中心（或 Hochschild 上同调）的谱之上，Hom 空间局部化为拟凝聚层，从而给出 $\mathcal C$ 的谱分解。因此，在研究表示的族——特别是中心特征变动的情形——很重要的场合，凝聚层会自然出现。在局部 Langlands 对应的背景下，我们仅举它与调和分析（例如见 [SV17]）、$K$-理论与 Baum--Connes 猜想 [ABPS17]、以及模表示与整表示论 [EH14,H20,HM18] 的联系作为众多动机中的几个。

Langlands 对应编码了这类雄心勃勃的代数谱分解，其中 Langlands 参数空间扮演中心参数的角色。一方面，（非分歧）Hecke 算子理论对约化群 $G$ 确定了自守形式（经典 Langlands 中）或自守层（几何 Langlands 中）上的一大交换对称代数。另一方面，这个交换代数的谱（经由 Satake 对应）与映入对偶群 $\check G$ 的 Langlands 参数空间相联系。于是我们希望用 Langlands 参数的代数几何来谱分解自守侧。这导向代数闭域上曲线的几何 Langlands 猜想（无论哪个变体）。\ysrc{425}该猜想断言：自守层的范畴可以等同于曲线上局部系统叠上（ind-）凝聚层的范畴。

近来出现了一种新的、受几何 Langlands 纲领启发的局部 Langlands 对应的「凝聚」范畴形式。其中，局部域上群的表示的范畴由 Langlands 参数叠上的凝聚层来描述（细节见本次暑期学校中 Fargues--Scholze、Emerton--Hellmann--Gee 与 Zhu 的讲义）。在阿基米德情形，这种凝聚表述在 [BN07, BN10b] 中作为扭量线上的驯顺几何 Langlands 对应被提出，本文也将讨论它。

在非阿基米德情形，这样一幅图景开始从若干不同方向浮现（特别见 [La18a, La18b, LaZ19, GL18, Ga16, GKRV20]），并导向 Fargues [Fa16]、Zhu [Zhu25] 与 Hellmann [He20]）的猜想\tn{底本 425 页此处多一个右括号。}。这一视角在 Fargues 与 Scholze [FS21] 的丰碑式工作中得到深刻发展；该工作尤其建立了「自守到伽罗瓦」的方向，证明自守范畴（因而是它里面由 $G(K)$ 的表示构成的子范畴）层化于 Langlands 参数叠之上。Zhu 即将发表的工作 [Zhu25] 运用范畴迹的理论（如 [Zh18] 中）与 Bezrukavnikov 的驯顺分歧局部几何 Langlands 对应 [Bez06, Bez16]，为幂幺表示建立凝聚的局部 Langlands 对应（其中主序列部分已在 [BCHN22] 中证明）。

叠上凝聚层的范畴比可构造层的范畴大得多。例如，在稳定化子群为 $H$ 的点处，拟凝聚层由 $H$ 的全部代数表示来标号，而可构造层只对应 $H$ 的分支群的表示（在导出情形则对应 $H$ 上链复形的模）。因此，为与 Langlands 参数上凝聚层的整个范畴相匹配，我们必须考虑一个大得多的表示论一侧（而不是某些特定的满子范畴）。事实上，[Zhu25,Fa16,FS21] 的猜想不是像 Vogan 猜想那样只涉及有限个纯内型，而是把 Langlands 参数上的凝聚层等同于无限多个群的表示。这些群都是 $G$ 的 Levi 因子的内型（的点群），作为 $G$-同晶的自同构群（或适当 $G$-丛叠中点的自同构群）出现，并由 Kottwitz 集 $B(G)$ 标号。

\subsection{代数与拓扑的联系}

粗略地说，我们描述了局部 Langlands 对应的以下两种风味：
\begin{center}
\begin{tikzpicture}[font=\small,
  box/.style={draw, inner sep=4pt, align=center, minimum height=1.6cm}]
\node[box, minimum width=5.4cm] (A) at (0,0) {$G(K)$ 的纯内型的表示\\（固定无穷小特征）};
\node[box, minimum width=4.8cm] (B) at (9.2cm,0) {Langlands 参数上的\\可构造层\\（固定无穷小特征）};
\node[box, minimum width=5.4cm] (C) at (0,-2.5cm) {$G$-同晶的稳定化子的\\光滑表示};
\node[box, minimum width=4.8cm] (D) at (9.2cm,-2.5cm) {Langlands 参数上的\\凝聚层};
\draw[dotted,->] (A.east) -- (B.west)
  node[midway,above=1pt]{\footnotesize「可构造」}
  node[midway,below=1pt]{\footnotesize 局部 Langlands};
\draw[dotted,->] (C.east) -- (D.west)
  node[midway,above=1pt]{\footnotesize「凝聚」}
  node[midway,below=1pt]{\footnotesize 局部 Langlands};
\foreach \n in {A,B,C,D}{
  \draw[thick] ([xshift=1.3pt,yshift=-1.3pt]\n.south west)
            -- ([xshift=1.3pt,yshift=-1.3pt]\n.south east)
            -- ([xshift=1.3pt,yshift=-1.3pt]\n.north east);
}
\end{tikzpicture}
\end{center}

\ysrc{426}
本讲的主要目标，是描述一个在局部 Langlands 纲领的凝聚设定与可构造设定之间插值的机制——具体地说，我们提出：后者作为前者某个满子范畴的形变中的一般点而出现。特别地，我们提出可构造范畴的一种形式，它确实随无穷小特征良好地变化；并提出一个（猜想性地）把表示论从全部 $G$-同晶砍到纯内型的操作。

设 $X$ 是光滑复簇。代数 de Rham 定理把常值层的上同调 $H^*(X,\C_X)\simeq\Ext^*_{\operatorname{Shv}(X)}(\C_X,\C_X)$——即该局部系统的导出自同态——等同于 de Rham 上同调 $H^*(X,\Omega^\bullet,d)\simeq\Ext^*_{\DD_X}(\OO_X,\OO_X)$——即带自然平坦联络的 $D$-模 $\OO_X$ 的导出自同态。Riemann--Hilbert 对应把它推广为 $X$ 上局部系统的阿贝尔范畴与 $X$ 上带正则平坦联络的向量丛的阿贝尔范畴之间的等价，并可延拓为 $X$ 的可构造导出范畴 $\operatorname{Shv}_c(X;\C)$ 与 $D_X$-模的导出范畴 $\mathcal D(X)$ 的一个满子范畴 $\mathcal D_{rh}(X)$（由正则完整（regular holonomic）对象组成）之间的范畴等价。这一等价经下降延拓到等变情形，即延拓到光滑叠 $X$ 上。此外，在几何表示论中最令人感兴趣的情形——即处理只有有限多个轨道的概形上的等变层——我们可以丢掉「正则完整」这一限定，并把等变可构造导出范畴等同于全部凝聚强等变 $D$-模的导出范畴 $\mathcal D_c(X)$。于是，例如在上面叙述的 Kazhdan--Lusztig 或 Springer 理论的情形，我们可以干脆用 $D$-模代替可构造层。

$D$-模本身可以描述为余切丛 $\mathbb T^*_X$ 上拟凝聚层的形变量子化。微分算子层 $\DD_X$ 按算子的阶过滤，其伴随分次 $\operatorname{gr}(\DD_X)\simeq\OO_{\mathbb T^*_X}$ 等同于符号，即 $X$ 的余切丛上的函数。因此 $\DD_X$ 可以描述为带标准辛形式的 $\OO_{\mathbb T^*_X}$ 的形变量子化：对 $\DD_X$ 作 Rees 构造，得到 $\mathbb A^1/\GG_m$ 上的一族 $\GG_m$-等变代数层 $\DD_{\hbar,X}$（$\hbar\in\mathbb A^1/\GG_m$），它在 $\hbar\neq0$ 处复原 $\DD_X$，在 $\hbar=0$ 处给出带 Poisson 括号的 $\OO_{\mathbb T^*_X}$，该括号来自交换子中的 $\hbar$-线性项。过渡到模之后，我们看到 $\DD_X$-模的范畴是 $\mathbb A^1$ 上一族 $\GG_m$-等变范畴在 $\hbar\neq0$ 处的纤维，并在 $\hbar=0$ 处退化到 $\QCoh(\mathbb T^*_X)^{\GG_m}$。

Connes 与 Feigin--Tsygan 的循环同调理论，经由导出代数几何的眼光来看，给出了 de Rham 复形——从而 $\DD$-模范畴——的拓扑解释。具体说，考虑单纯（或「动态的」，animated）圆周 $S^1=B\Z$ 与导出圈空间 $\mathcal L X=\Map(S^1,X)=X\times_{X\times X}X$。在概形上，Hochschild--Kostant--Rosenberg 定理把 $\mathcal L X$ 等同于微分形式代数（置于负上同调次数）的谱；而 de Rham 微分则有令人愉快的解释：它是 $\mathcal L X$ 上圈旋转 $S^1$-作用的线性化，并由次数为 $-1$ 的 Connes $B$-算子编码。\ysrc{427}Rees 参数 $\hbar\in\mathbb A^1$ 被解释为（差一个上同调平移）$S^1$-等变参数 $k[u]\simeq H^*_{S^1}(\mathrm{pt})\simeq H^*(\C\PP^\infty)$。于是（差上同调平移与重整化）从 $\QCoh(\mathbb T^*_X)$ 到 $X$ 上 $D$-模的 $\hbar$-形变，等同于 $\mathcal L X$ 上\emph{循环层}（cyclic sheaves）的 $k[u]$-线性范畴 $\QCoh^!(\mathcal L X)^{S^1}$，即圈空间上的 $S^1$-等变（ind-）凝聚层，也就是循环同调理论的范畴化。过渡到\emph{周期循环层}——即把等变参数 $u$ 取逆——我们复原 $\DD_X$-模范畴的基变换 $\DD_X\text{-mod}\otimes_k k[u,u^{-1}]$。

在光滑概形的情形，这一关于 $D$-模的循环视角可以看作形式操作；但在光滑叠 $X$ 的情形，它的味道就完全不同了。导出圈空间 $\mathcal L X$ 等同于 $X$ 的（导出）惯性叠，即点与自同构配对的叠
\[
\mathcal L X\simeq\{x\in X,\ g\in\Aut(x)\};
\]
在有限轨道叠这一特殊情形，它实际上是普通的（非导出的）叠。上述 Koszul 对偶模式的一个推广，把光滑叠 $X$ 上 $D$-模的范畴与圈空间形式完备化上的圆作用联系起来（这里把 $\Aut(x)$ 换成它的形式群）。但 $\mathcal L X$ 上凝聚层的整个范畴（连同其圆作用）给出一个更新、更丰富的范畴，它通过添加新的连续参数（「中心特征」，表现为圈 $g$ 的特征值）来增强 $X$ 上的 $D$-模。

本讲的要旨是：这类圈空间及其变体作为 Langlands 参数叠自然出现，它们的凝聚层范畴为局部 Langlands 对应的范畴形式提供了自然的谱侧。关键在于，这一描述给这些范畴配备了圆作用。在阿基米德情形，圆作用来自 $S^1$ 在扭量线上的几何作用。在幂幺的非阿基米德情形，圆作用来自通过变动 $q$ 在 Frobenius 的迹与恒等同态（圈空间）的迹之间插值（非正式地说，它旋转「圆周」$\Spec(\F_q)$，其基本群由 Frobenius 生成）。一般情形下，我们只是猜测存在圆作用的一般构造，它反映 Frobenius 迹理论中某个微妙之处。

圆作用的用处在于：它们提供了所需的机制，把凝聚层与固定中心特征的 Langlands 参数空间上的 $D$-模联系起来，即在凝聚与可构造两种 Langlands 对应之间插值。也就是说，循环层（即 $S^1$-等变层）的范畴构成等变参数 $u$ 上的一族。取 $u=0$，我们得到（$S^1$-不变的）凝聚层的一个满子范畴，即「几何 Langlands 型」局部 Langlands 对应的设定。另一方面，对非零的 $u$，我们得到那些经由 Vogan 局部 Langlands 对应来参数化纯内型表示的 $D$-模范畴。

\subsection{Jordan 分解}

联系循环层与表示论的关键机制是「圈上的 Jordan 分解」：它是约化群中共轭类的 Jordan 分解的叠化推广，使我们能在不同中心特征之间插值。让我们来说明最基本的例子。

考虑叠 $\check G/\check G$（回忆它是叠 $B\check G=\mathrm{pt}/\check G$ 中的惯性，即导出圈 $\check G/\check G=\mathcal L(B\check G)$）。\ysrc{428}取不变多项式（即过渡到仿射化）定义映射
\[
\chi:\ \check G/\check G\longrightarrow\check G\sslash\check G=\Spec(\OO(\check G)^{\check G})\simeq\check H\sslash W
\]
映到连续参数的空间。若像局部 Langlands 对应的可构造形式那样，把 $\check G\sslash\check G$ 解释为 $\check G$ 中的半单共轭类，我们就会遇到一个明显的不连续性：取定类 $[\alpha]\in\check G\sslash\check G$ 的一个半单代表元 $\alpha\in(\check G)^{\mathrm{ss}}$，相应的中心化子 $\check G(\alpha)$ 会不连续地依赖于特征值；同样，分类叠 $B\check G(\alpha)$ 在 $\check H\sslash W$ 上并不构成好的族。然而 $\check G$ 中的 Jordan 分解通过添加交换的幂幺元修正了这一点：$[\alpha]$ 的形式邻域的纤维 $\chi^{-1}(\widehat{[\alpha]})$ 等同于 $\check G(\alpha)^u/\check G(\alpha)$，即中心化子 $\check G(\alpha)$ 中幂幺类的叠。这是更一般构造——幂幺圈空间——的特例：它是惯性对任意叠的变体，其中只允许幂幺自同构：
\[
\mathcal L^u(B\check G(\alpha))=\check G(\alpha)^u/\check G(\alpha).
\]
换言之，若把 $B\check G(\alpha)$ 扩充为允许 $\check G(\alpha)$ 中的非平凡幂幺类（即用它替换为它的幂幺圈），我们就得到一个好的族，其全空间是 $B\check G$ 中的圈。

幂幺自同构（或抽象地说幂幺圈）的一个关键特征是：它们带有把自身收缩到恒同的 $\GG_m$ 作用。这一作用使我们能把全纤维 $\chi^{-1}(\widehat{[\alpha]})=\check G(\alpha)^u/\check G(\alpha)$ 上的凝聚层与其形式完备化上的凝聚层联系起来。最后，Koszul 对偶把这一形式完备化上的循环层与叠 $B\check G(\alpha)$ 本身上的 $D$-模联系起来。于是，$B\check G(\alpha)$ 上 $D$-模范畴的族，可以用 $\check G/\check G$ 上凝聚层范畴上的圆作用统一地表达。

我们将在阿基米德与非阿基米德两种设定中看到类似现象：放弃要求 Frobenius 或阿基米德 Weil 群的作用半单，就得到 Langlands 参数空间的若干版本，它们光滑地在不同中心特征之间插值，却由于收缩 $\GG_m$ 作用与 $S^1$-等变局部化的结合，仍然捕捉到相同的可构造层范畴。此外，这些更大的 Langlands 参数空间上的凝聚层范畴，恰好就是局部 Langlands 对应的几何 Langlands 风格表述所需要的。

\subsection{全文结构}

本文余下内容的组织如下：第 2 节是对导出圈空间的简短而平易的介绍。在第 3、4 节，我们分别从圈空间与圆作用的视角描述阿基米德与非阿基米德局部 Langlands 对应，而把导出代数几何中底层技术机制的概述推到 A 节。更详细地说：

\ysrc{429}在第 3 节，我们描述圈空间在幂幺非阿基米德局部 Langlands 对应中的作用。我们首先描述范畴局部 Langlands 纲领中的一些近期结果，这些结果把表示的范畴实现为凝聚层的满子范畴。然后描述我们的预印本 [BCHN] 中发展的结果与猜想，它们涉及圆作用在幂幺局部 Langlands 对应中作为凝聚形式与可构造形式之间的插值所扮演的角色。随后我们简要猜测圆作用在一般非阿基米德局部 Langlands 对应中的出现。

在第 4 节，我们描述圈空间在阿基米德局部 Langlands 对应中的作用。特别地，我们按 [BN10b] 给出 ABV 几何参数空间 [ABV92] 的一个简单叠论描述，以及 Soergel 猜想的一个相应陈述。然后我们引入与扭量线 $\widetilde{\mathbb P}^1_\R$ 相伴的 Langlands 参数叠；它在无穷小特征变动时于 ABV 参数空间之间光滑插值。我们并把这个设定下的 Soergel 猜想推广为用循环层表述的猜想。在自守侧，我们从扭量线上的实抛物丛叠复原纯内型的表示。我们（按 [BN07,BN16]）对 $\widetilde{\mathbb P}^1_\R$ 表述驯顺分歧几何 Langlands 猜想，并解释它如何以周期循环形变的形式复原 Soergel 猜想。

在附录 A 节，我们描述本文赖以成立的关键技术机制。我们分三步把圈空间（其中包括 Langlands 参数叠）上 $S^1$-等变凝聚层的范畴与 $D$-模（特别是在 Vogan 簇上）的范畴联系起来。第一步是圈上的 Jordan 分解——一个把商叠的圈空间与幂幺圈空间联系起来的普遍模式。其次，幂幺圈空间上的收缩 $\GG_m$ 作用把它们的层论与形式圈的层论联系起来。最后，我们应用形式圈空间上 $S^1$-等变凝聚层与 $D$-模之间的 Koszul 对偶。

\section{圈空间上的等变层}

本节小结导出代数几何中圈空间上等变层的关键陈述；这正是后面各节表示论构造所处的设定。我们所处的同伦代数一般设定的参考书是 [Lur1, Lur2]；不过我们会尽量压制 $\infty$-范畴的语言。$S^1$-作用在导出圈空间背景下的讨论见 [TV]。我们的目标是为其中一些构造给出可用的直观；一般理论的进一步细节我们推迟到 A 节。

我们采用一个适用面很广的设定：固定特征零的域 $k$，取具有仿射对角线的光滑 Artin 叠 $X$。在所有当前感兴趣的情形，都可假设 $X=Y/G$，其中 $Y$ 是光滑拟射影概形，仿射代数群 $G$ 作用于其上。

\subsection{圈的各种变体}

我们把拓扑群 $S^1=B\Z=K(\Z,1)$ 看作预叠中的群对象，即从环（或概形）到拓扑空间（或单纯集）同伦论中的群的［局部］常值函子。\ysrc{430}我们可以把它看作 $S^0=\Spec k\sqcup\Spec k$ 的悬置：
\[
S^1\simeq\Sigma S^0=\Spec k\coprod_{\Spec k\sqcup\Spec k}\Spec k.
\]
$S^1$ 作为余极限的这一表示，引出从 $S^1$ 出发的映射叠——即（导出）\emph{圈空间}——作为极限的一个表示：
\[
\mathcal L X:=\Map(S^1,X)=\Map\Bigl(\Spec k\coprod_{\Spec k\sqcup\Spec k}\Spec k,\ X\Bigr)=X\times_{X\times X}X.
\]
这里（导出）纤维积是 $X$ 的对角线的自交。等价地说，圈空间是（导出）惯性叠，其 $k$-点 $\mathcal L X(k)$ 由如下的对 $(x,\gamma)$ 给出（模去 $\Aut(x)$-共轭）：$x\in X(k)$，而 $\gamma\in\Aut(x)=(\{x\}\times_X\{x\})(k)$ 是点 $x$ 的一个自同构。

\begin{exam}{2.1}\label{exam:a5-01-1}
我们考虑以下例子。
\begin{enumerate}[(1),itemsep=2pt]
\item 设 $X$ 是概形。记 $\mathcal B_X$ 为 $X$ 上层化的循环 bar 复形，把它看作带 shuffle 积的 dg 代数层；则
\[
\mathcal L X=\Spec_X\mathcal B_X.
\]
\item 设 $X=BG$ 是分类叠。此时对角线光滑，因而平坦。于是 $\mathcal L X$ 是非导出叠，即经典惯性叠 $G/G$。
\item 设 $X=Y/G$ 是整体商叠。圈空间有如下表示，即沿可表示映射的叠的导出纤维积：
\begin{center}
\begin{tikzcd}
\mathcal L(Y/G) \arrow[r] \arrow[d] & (Y\times G)/G \arrow[d, "{(a,p)}"] \\
Y/G \arrow[r, "{\Delta}"] & (Y\times Y)/G
\end{tikzcd}
\end{center}
其中 $G$ 对角地作用于 $Y\times Y$ 与 $Y\times G$，$a$ 是作用映射，$p$ 是投影映射，$\Delta$ 是对角映射。这个叠的 $k$-点为
\[
\mathcal L(Y/G)(k)=\{(y,g)\in Y(k)\times G(k)\mid g\cdot y=y\}/G(k).
\]
\item 设 $X$ 是光滑有限轨道叠（即只有有限多个 $k$-点的叠）。此时导出圈空间 $\mathcal L X$ 实际上是普通的（非导出）叠，并等同于 $X$ 的经典惯性叠。这种导出结构的缺失并不依赖于例 2.1 中图里任何映射的平坦性，而是一个维数计数：若 $X$ 是有限轨道叠，则 $\mathcal L X$ 是若干（叠维数为 $0$ 的）不可约分支的并，而这个维数等于预期维数。\footnote{虽然这是众所周知的，但由于它很基本，我们为不熟悉的读者概述一下论证。首先，任何纤维积 $X\times_YZ$ 都可以写成 $X\times Z$ 与 $X\times Z$（在不同映射下）在 $X\times Y\times Z$ 中的交，因此可设我们的导出纤维积是导出交 $X\cap_YZ$。局部地，由于一切光滑，$X,Z$ 由正则序列切出；又因 $\operatorname{codim}(X\cap_YZ)=\operatorname{codim}(X)+\operatorname{codim}(Z)$，这些序列的并也是正则的，于是高阶 $\Tor$ 消失。}
\end{enumerate}
\end{exam}

\subsection{圈旋转}

圈空间 $\mathcal L X=\Map(S^1,X)$ 从拓扑继承了许多熟悉的代数结构，我们关注其中以下三个：
\begin{enumerate}[(1),itemsep=2pt]
\item 限制到点 $\mathrm{pt}\in S^1$ 给出到基点的投影 $\mathcal L X\to X$。
\item 沿 $S^1\to\mathrm{pt}$ 的拉回给出常值圈 $X\subset\mathcal L X$。
\item $S^1$ 通过旋转作用于自身，给出 $\mathcal L X$ 上由圈旋转而来的 $S^1$-作用（见 A.1 节）。
\end{enumerate}

\ysrc{431}它在我们的展开中不起作用，但也值得提及：坍缩映射 $S^1\to S^1\vee S^1$、逆映射 $S^1\to S^1$ 以及它们的相容性，给出 $\mathcal L X\to X$ 上的\emph{相对群结构}。

\subsubsection{什么是圆作用？}

我们的主要目标之一，是把圈旋转作用具体地阐明。更一般地，应当如何理解导出代数几何中的 $S^1$-作用？（关于圆作用的一般同伦论讨论，例如见 [Lur3]。）换言之，导出叠中群对象的映射 $S^1\to\Aut(Z)$ 究竟相当于什么？对这一结构的精确描述需要动用 $\infty$-范畴，用它来以精简的方式组织同伦相容的结构。

把 $S^1=B\Z$ 等同之后，我们可以忘掉群结构，考虑底层的叠映射 $B\Z\to\Aut(Z)$，它相当于 $Z$ 的恒等同态的一个自同构，即每个点上的一个典范自同构。然而，仅给出这一映射并不足以完整描述一个 $S^1$-作用。思考叠 $Z$ 上 $S^1$-作用的另一种简洁方式是这样一组数据：一个叠 $W$，带有到 $BS^1\simeq K(\Z,2)$ 的投影（$BS^1$ 是对应于 $\C\PP^\infty$ 同伦型的叠），以及一个等价 $W\times_{BS^1}\mathrm{pt}\simeq Z$。这里，仅给出作用映射这一数据，对应于把该投影限制到 $2$-球面 $\C\PP^1\subset\C\PP^\infty$；在这一表述中，高阶相容性相当于把该映射逐次延拓到更高维的胞腔上。

通过线性化，可以得到圆作用的其他具体模型。在函数链复形的层面，我们可以线性化 $S^1$ 作用，得到 $S^1$ 的上链代数（或上同调代数\footnote{提醒：我们只在特征 $0$ 下工作！}）这一 Hopf 代数在 $Z$ 上（导出）交换函数环上的作用，它相当于一个上同调次数为 $-1$ 的平方零导子。在层范畴的层面，可以用 $\Z$ 与 $\GG_m$ 之间的 Cartier 对偶（即 Fourier 级数）把带 $S^1$-作用的（稳定）范畴等同于 $\GG_m$-线性范畴。

直观上，我们把圈空间上的圆作用看作两种简单现象的组合，它们分别对应圆作用作为恒同的自同构的化身与作为平方零奇导子的化身。在一个极端，当 $X=BG$（$G$ 有限）是有限轨形时，$\mathcal L X=G/G$（伴随商）上的圆作用（通过恒同的自同构来看）归结为一句关于中心化子的非常具体的陈述，见下面的例 2.2；而按圈旋转作等变工作，意味着考虑中心化子的、在中心元素上作用平凡的表示。在另一个极端，当 $X$ 是概形时，它（通过线性化）归结为 de Rham 微分的一种替代的「拓扑」观点；而按圈旋转作等变工作，意味着考虑带平坦联络的层。一般情形下，很难完整描述圈旋转圆作用。然而我们将说明：更仔细地考察圈空间的几何，就能把圆作用的研究（就等变层而言）归结为这两种现象的组合。详细讨论与更多动手的例子见附录。

\begin{exam}{2.2}\label{exam:a5-01-2}
设 $X=BG$，$G$ 有限。此时可以把伴随商 $\mathcal L X=G/G$ 等同于分类叠 $\coprod_{[\alpha]}BG(\alpha)$ 的不交并，其中 $G(\alpha)\subset G$ 是群元素 $\alpha$ 的中心化子，$\alpha$ 取遍共轭类 $[\alpha]$。注意中心化子 $G(\alpha)$ 的中心有一个典范元素，即 $\alpha$ 本身；\ysrc{432}于是 $\mathcal L X$ 的每个点都带有一个典范自同构，这正是圈旋转 $S^1$ 作用的一种体现。

换一种说法：对每个分支 $BG(\alpha)$，其自同构空间是同伦商 $\Aut(G(\alpha))/G$，其中 $G$ 通过复合以共轭作用；恒同自同构的稳定化子是 $Z(G(\alpha))$。于是分支 $BG(\alpha)$ 上的 $S^1$-作用由中心元素 $\alpha\in Z(G(\alpha))$ 给出；把它看作 $BG(\alpha)$ 恒同映射的自同构，即任何 $G(\alpha)$-丛的一个典范自同构。
\end{exam}

\subsubsection{奇切丛}

出于典型的原因，我们可能试图通过线性化来理解复杂的非线性对象。导出圈空间相应的线性化是\emph{奇切丛}
\[
\mathbb T_X[-1]=\Spec_X\Sym^\bullet_X\Omega^1_X,
\]
其中 $\Omega^1_X$ 是 $X$ 的余切复形，即导出的 Kähler 微分。奇切丛可以等同于常值圈 $X\subset\mathcal L X$ 的法丛。它是 $X$ 上的导出向量丛，因而带有自然的收缩 $\GG_m$ 作用。

对 $S^1=B\Z$ 自身的线性化，我们可以取其仿射化 $B\GG_a=\Spec\OO(S^1)$，即连通 $k$-代数上的一个函子，它由 $S^1$ 上的函数 $\OO(S^1)\simeq C^\bullet(S^1;k)$ 表示，也就是由 $S^1$ 的 $k$-值上链表示。\footnote{注意，由于 $\OO(S^1)$ 是余连通的（而非连通的），它的「谱」不是仿射导出概形，而是 Toën 意义下的仿射叠（或 Lurie 意义下的余仿射叠）。}群 $\GG_a$ 通过诱导 $B\GG_a$ 上收缩 $\GG_m$-作用的群同态，带有 $\GG_m$ 的一个典范收缩作用。用更接地气的话说，$\OO(S^1)\simeq C^\bullet(S^1;k)$ 作为 dg 代数是形式的，因而带有与上同调次数一致的一个典范内蕴权分次。

我们还会对\emph{形式奇切丛} $\widehat{\mathbb T}_X[-1]$ 感兴趣，即 $\mathbb T_X[-1]$ 沿其零截面的形式完备化。$\mathcal L X$ 上的 $S^1$-作用诱导 $\mathbb T_X[-1]$ 与 $\widehat{\mathbb T}_X[-1]$ 上的 $S^1$-作用，而后者通过仿射化映射 $S^1\to B\GG_a$ 分解 [Ch23, Prop. 3.2.3]。我们可以用一个半直积 $B\GG_a\rtimes\GG_m$ 的作用来统一整理 $\widehat{\mathbb T}_X[-1]$ 上的这些作用。

\begin{exam}{2.3}\label{exam:a5-01-3}
我们继续例 2.1 中的例子。
\begin{enumerate}[(1),itemsep=2pt]
\item 设 $X$ 是概形。Hochschild--Kostant--Rosenberg（或指数）映射定义一个 $S^1$-等变（且通过 $B\GG_a$-等变分解）的等价（见 [BN10a, 命题 4.4]）：
\[
\exp:\ \mathbb T_X[-1]\xrightarrow{\ \sim\ }\mathcal L X.
\]
特别地，当 $X$ 是概形时，$\Omega^1_X[1]$ 具有负的 Tor 幅度，因此 $X\subset\mathbb T_X[-1]$ 是幂零加厚。于是 $\mathcal L X\simeq\mathbb T_X[-1]\simeq\widehat{\mathbb T}_X[-1]$。所以对概形 $X$，导出圈空间 $\mathcal L X$ 并不比它的线性化更复杂。
\item 设 $X=BG$ 是分类叠。奇切丛 $\mathbb T_X[-1]=\mathfrak g/G$ 是伴随等变的 Lie 代数，形式奇切丛 $\widehat{\mathbb T}_X[-1]=\widehat{\mathfrak g}/G$ 是伴随等变的形式 Lie 代数。
\end{enumerate}
\end{exam}
